\subsection{Müste'menlere Uygulanan Ceza Hukuku}
\indent\indent İslam ülkesi vatandaşı olmayıp, devlet tarafından verilen geçici süre izinle İslam ükesinde bulunan kişiler müste'men olarak adlandırılıyorlardı. 
Bu kimselerin İslam ve Osmanlı hukuku içindeki statüleri büyük ölçüde zımmîlere benziyordu. Hanefî mezhebine göre, şahıs haklarını ihlal eden suçlarda 
müste'menlere İslam hukuku uygulanırdı. Bu bağlamda müslüman ya da zımmîyi öldüren müste'mene, şartlara bağlı olarak kısas cezası uygulanırdı. Zina ve 
içki içme gibi suçlarda ise İslam hukuku uygulanmazdı. Öte yenadan Şâfiî, Mâlikî ve Hanbelî mezhepleri Hanefî  mezhebinden farklı olarak, her türlü suçta 
müste'menlere İslam hukukunun uygulanmasını savunuyorlardı.\footcite[324]{osmanli_hukuku}

16.yüzyılda uygulamaya konulan kapitülasyonlar ile yabancı devletler hukuki alanda da bazı imtiyazlar elde etmişlerdir. Fransa'ya 1535 yılında verilen 
kapitülasyonlar ile Fransa devletine tabiî tüccarların kendi aralarındaki anlaşmazlıklarda, bu devletin konsoloslukları yetkili kılınmıştır. Kapitülasyonlar,
zaman içinde diğer devletlerin de faydalanacağı şekilde genişletilmiştir. Böylece, müste'menler giderek Osmanlı mahkemelerinin yargı alanının dışında kalmaya 
başlamışlardır. Müste'menlerin Osmanlı teb`asına karşı işlediği suçlarda Osmanlı mahkemeleri yetkili olmaya devam etse de, bazı davaların Dîvân-ı Hûmâyun'da 
görülmesi ve yargılama esnasında bir elçilik görevlisinin hazır bulunması gibi bir takım ayrıcalıklar tanınmıştır.\footcite[325]{osmanli_hukuku}